\chapter{Introduzione}

In questo capitolo introduttivo della tesi viene delineato uno dei temi più
rilevanti legati all'utilizzo dell'intelligenza artificiale (IA), ovvero
l'impatto energetico e ambientale dei \textit{Large Language Model (LLM)}.
La tesi si inserisce nell'ambito della ricerca sull'IA sostenibile,
concentrandosi sulla sperimentazione di una delle sue frontiere emergenti:
il \textit{Green Edge AI}. Successivamente, vengono descritti gli obiettivi
principali del lavoro.

\section{Background e Motivazioni}

Negli ultimi anni, l'utilizzo dei \textit{LLM} ha conosciuto una rapida
diffusione, trovando applicazione in numerosi ambiti e scenari operativi.
Questa crescita ha comportato un significativo incremento dei consumi
energetici e ha imposto requisiti sempre più elevati alle infrastrutture
di rete incaricate di gestire il crescente carico computazionale e di
traffico.

Un approccio innovativo proposto dalla ricerca consiste nello spostamento
delle capacità di elaborazione intelligente dalle infrastrutture cloud verso
dispositivi \textit{edge} e \textit{Internet of Things (IoT)}, avvicinando
l'esecuzione dei modelli al dispositivo dell'utente finale. Tale paradigma,
oltre a ridurre la dipendenza dalle infrastrutture cloud centralizzate,
consente una maggiore tutela della privacy dei dati e una riduzione della
latenza delle applicazioni.

Questa evoluzione è resa possibile dallo sviluppo di tecniche di
quantizzazione quali la \textit{Post-Training Quantization (PTQ)} e la
\textit{Quantization-Aware Training (QAT)}, che permettono di ridurre la
precisione numerica con cui vengono rappresentati i pesi delle reti neurali.
In questo modo, dispositivi \textit{edge} e \textit{IoT}, caratterizzati da
risorse computazionali e di memoria limitate, possono supportare l'esecuzione
locale di LLM. Tuttavia, tali tecniche introducono generalmente una
degradazione della qualità delle risposte generate dai modelli e non
garantiscono necessariamente un miglioramento ottimale delle prestazioni
di latenza.

\section{Obiettivi della Tesi}

La presente tesi propone un benchmark relativo a tre aspetti fondamentali
dell'esecuzione locale di LLM quantizzati: \textit{efficienza energetica},
\textit{qualità delle risposte} e \textit{latenza di inferenza}. La valutazione
viene effettuata su sei LLM quantizzati mediante tecnica \textit{Post-Training
Quantization (PTQ)}, eseguiti localmente su un dispositivo
\textbf{Raspberry Pi 5 con 8 GB di RAM}.

I modelli analizzati sono riportati nella Tabella~\ref{tab:llm}. Come si può
osservare, sono state considerate tre differenti famiglie di modelli, per
ciascuna delle quali sono stati valutati due livelli di quantizzazione:
\textit{4-bit} e \textit{8-bit}. Inoltre, i modelli sono stati selezionati
considerando tre differenti dimensioni in termini di numero di parametri:
1.5 miliardi, 2 miliardi e 3 miliardi. In questo modo, la tesi fornisce una
valutazione rappresentativa di una parte della gamma di dimensioni dei modelli
LLM eseguibili su Raspberry Pi 5.

\begin{table}[htbp]
    \centering
    \begin{tabular}{lcc}
        \hline
        Modello & Parametri & Livello di quantizzazione \\
        \hline
        Qwen2.5 & 1.5B & 4 bit \\
        Qwen2.5 & 1.5B & 8 bit \\
        Gemma2 & 2B & 4 bit \\
        Gemma2 & 2B & 8 bit \\
        Llama3.2 & 3B & 4 bit \\
        Llama3.2 & 3B & 8 bit \\
        \hline
    \end{tabular}
    \caption{\textit{LLM} quantizzati utilizzati nell'esperimento di tesi.}
    \label{tab:llm}
\end{table}

Per ciascuno dei sei modelli vengono analizzati tre indicatori principali:
l'\textit{efficienza energetica}, la \textit{qualità della risposta} generata e la \textit{latenza di
inferenza}. La valutazione viene condotta mediante cinque differenti benchmark
di inferenza, ciascuno rappresentato da un dataset standardizzato selezionato
per analizzare specifiche capacità degli LLM in scenari caratterizzati da
risorse computazionali limitate.

I benchmark utilizzati sono:

\begin{itemize}
    \item \textit{CommonsenseQA}: benchmark finalizzato alla valutazione delle
    capacità di ragionamento basate sul senso comune attraverso quesiti a
    scelta multipla;

    \item \textit{BIG-Bench Hard}: insieme di task complessi derivati dal
    benchmark BIG-Bench, utilizzato per misurare capacità avanzate di
    ragionamento e generalizzazione;

    \item \textit{TruthfulQA}: benchmark progettato per valutare la
    correttezza e l'affidabilità delle risposte generate dal modello in
    presenza di domande soggette a misconcezioni o informazioni errate;

    \item \textit{GSM8K}: benchmark composto da problemi matematici elementari
    che richiedono ragionamento aritmetico articolato in più passaggi;

    \item \textit{HumanEval}: benchmark per la valutazione della generazione
    automatica di codice tramite esercizi di programmazione verificati
    attraverso apposite suite di test.
\end{itemize}

Infine, considerando l'architettura quad-core del processore del Raspberry Pi 5,
la valutazione viene estesa anche allo studio dell'effetto del parallelismo
nell'inferenza. Per ciascun modello vengono quindi analizzate tre configurazioni
relative al numero di thread di esecuzione (1, 2 e 4 thread), con l'obiettivo
di studiare la relazione tra grado di parallelismo, consumo energetico e
latenza di inferenza.

Nel capitolo successivo viene descritto il setup sperimentale adottato per
l'acquisizione delle misure considerate nell'analisi.
